\documentclass[a4paper,12pt]{article}
\usepackage[utf8]{inputenc}
\usepackage[T1]{fontenc}
\usepackage[ngerman]{babel}
\usepackage{amsmath}
\usepackage{amssymb}
\usepackage{amsthm}
\usepackage{enumitem}
\usepackage{geometry}
\geometry{a4paper, margin=2.5cm}

\title{Einsendeaufgaben -- Lektion 1\\[0.5em]
\large Modul 61112: Lineare Algebra}
\author{}
\date{}

\begin{document}
\maketitle

\section*{Aufgabe 1.2}

\begin{enumerate}[label=\alph*)]

\item Sei $x \in G$ beliebig. Dann erfüllt $G_x = \{g \in G \mid x = g^{-1} \circ x \circ g\}$
die Bedingungen einer Untergruppe:

\textbf{(UG1)} Das neutrale Element von $G$ liegt in $G_x$, denn
\[
  e^{-1} \circ x \circ e = e^{-1} \circ x = x
\]

\textbf{(UG2)} $G_x$ ist abgeschlossen, denn für alle $a, b \in G_x$ gilt:
\begin{align*}
  x &= b^{-1} \circ x \circ b
    && \text{(wegen } b \in G_x\text{)} \\
    &= b^{-1} \circ (a^{-1} \circ x \circ a) \circ b
    && \text{(wegen } a \in G_x\text{)} \\
    &= (b^{-1} \circ a^{-1}) \circ x \circ (a \circ b)
    && \text{(wegen Assoziativität in } G\text{)} \\
    &= (a \circ b)^{-1} \circ x \circ (a \circ b)
    && \text{(wegen } (a \circ b) \circ (b^{-1} \circ a^{-1}) = e\text{)}
\end{align*}

\textbf{(UG3)} Für alle $g \in G_x$ gilt $g^{-1} \in G_x$, denn
\begin{align*}
  x &= (g \circ g^{-1}) \circ x \circ (g \circ g^{-1})
    && \text{(wegen } g \circ g^{-1} = e\text{)} \\
    &= g \circ (g^{-1} \circ x \circ g) \circ g^{-1}
    && \text{(wegen Assoziativität in } G\text{)} \\
    &= g \circ x \circ g^{-1}
    && \text{(wegen } g \in G_x\text{)}
\end{align*}

\item Sei $x, y \in G$ beliebig. Wir zeigen, dass für alle $g \in G_x \cap G_y$ gilt
$g \in G_{x \circ y}$. Also sei $g \in G_x \cap G_y$ beliebig.

Es gilt
\begin{align*}
  &g^{-1} \circ (x \circ y) \circ g \\
  &= g^{-1} \circ \bigl((g \circ x \circ g^{-1}) \circ (g \circ y \circ g^{-1})\bigr) \circ g
    && \text{(wegen } g \in G_x,\; g \in G_y \\
  & && \text{und UG3} \\
  &= (g^{-1} \circ g) \circ x \circ (g^{-1} \circ g) \circ y \circ (g^{-1} \circ g)
    && \text{(wegen Assoziativität in } G\text{)} \\
  &= x \circ y
\end{align*}

\end{enumerate}

\end{document}
